% paper2.tex — Topological phonon gaps in Fibonacci isotope nanoribbons of h-BN
% Compile: pdflatex paper2 && bibtex paper2 && pdflatex paper2 && pdflatex paper2
\documentclass[aps,prl,preprint,superscriptaddress,longbibliography]{revtex4-2}
% switch 'preprint' -> 'reprint' for the journal two-column layout
\usepackage{graphicx}
\usepackage{amsmath,amssymb,bm}
\usepackage[colorlinks=true,linkcolor=blue,citecolor=blue,urlcolor=blue]{hyperref}

\graphicspath{{figures/}}
\begin{document}

\title{Topological phonon gaps in Fibonacci isotope nanoribbons of hexagonal boron nitride}

\author{Ronghua Wu}
\email{wrh@gdtzn.com}
\affiliation{Zhuhai Jiuchongtian Aviation Technology Co., Ltd., Zhuhai, China}
\author{Ruqing Chen}
\email{ruqing@hotmail.com}
\affiliation{GUT Geoservice Inc., Montreal, Canada}

\date{\today}

\begin{abstract}
Hyperuniform isotope order suppresses phonon scattering through a power law
that transverse channel averaging cannot disturb [R.~Wu and R.~Chen, Zenodo
(2026)]. Here we invert the objective and ask for the opposite spectral
object: a topologically labelled phonon gap, which the same averaging
destroys. Survival therefore demands a strictly single-mode nanoribbon, and
with it a confrontation with edge-roughness decoherence. Anchoring the
roughness scattering to a grid-converged boundary-conforming Green's-function
calculation, validated against an analytic Born result, we find that the
Ziman ray formula underestimates the coherence length by an order of
magnitude in the single-mode wave limit $\lambda \gg W$: a zero-parameter law
$\xi_{\mathrm{edge}} = 2W^2/[q^2 S_w(2q)]$ reproduces a 15-point numerical
scan with all three exponents (frequency to $10^{-3}$) and magnitude to $3.5\%$, softening the edge-flatness requirement
from atomic ($\lesssim 0.13$~nm) to lithographic ($\lesssim 1$~nm). A
Fibonacci $^{10}$BN/$^{11}$BN modulation of period $24.8$~nm in a $30$~nm
ribbon places a Chern-labelled gap at $0.250$~THz at a factor $2.20$ cost in
localization length relative to a periodic superlattice. Gap labelling is
verified falsifiably (median residual $3\times10^{-6}$, $1385\times$ below
random controls) and phason pumping exhibits the full counter-propagating
edge-state structure, with per-edge spectral flow $\pm q$ for all twelve
principal gaps up to $|q| = 11$. We propose an asynchronous-optical-sampling
cavity measurement with a built-in opposite-Chern-number control gap.
\end{abstract}

\maketitle

\section{Introduction}

In a companion work~\cite{paper1} we established that the relaxation rate for
phonon scattering by a spatially correlated isotope arrangement obeys
$\tau^{-1}_{\rm iso} \propto \omega^{\,d+1+\alpha}$, the hyperuniformity
exponent $\alpha$ of the arrangement~\cite{torquato2003} entering as an
additive, dimension-independent correction to the Rayleigh--Tamura
exponent~\cite{tamura1983}. The conclusion drawn there was suppressive:
ordering the isotopes \emph{removes} long-wavelength scattering. The present
work inverts the objective, and in doing so inverts every geometric
conclusion that followed from it.

\emph{The mirror.}---A power law is a scale invariant. Averaging it over the
transverse channels of a three-dimensional sample changes the prefactor and
leaves the exponent untouched; that immunity licensed the thick-sample
geometry of Ref.~\cite{paper1}, and made the nanoribbon actively harmful
there. Here we seek a gap. A gap is not scale invariant: its position in
transverse channel $m$ is set by the channel's own longitudinal wavevector,
$q_m^2 = (\omega^2 - v^2 k_m^2)/v^2$, so a channel sum migrates the gap
frequency channel by channel and smears a Cantor set of gaps into a
featureless continuum. The very averaging that protects a power law destroys
a gap. Transverse quantisation is therefore not a preference but a
prerequisite, in its strictest form: single-mode operation,
$N_{\rm prop} = 1$. That admits only the one-dimensional nanoribbon --- the
geometry Ref.~\cite{paper1} discarded --- and hands us the edge-decoherence
problem it avoided. Section~\ref{sec:survival} is that confrontation. The two
papers do not disagree: they report one physical fact, that channel averaging
preserves exponents and destroys gaps, read under opposite objectives.

\emph{Why Fibonacci, and not the primes.}---A gap alone is not topological.
The topological content of a one-dimensional quasiperiodic chain rests on the
gap-labelling theorem~\cite{bellissard1992}, which applies to sequences with
a pure-point diffraction spectrum and a cut-and-project construction:
Fibonacci~\cite{kohmoto1983,ostlund1983}, Sturmian words, the quadratic
irrationals. Under cut-and-project the chain is a slice of a two-dimensional
periodic parent, the slice offset $\phi$ is a synthetic second
dimension~\cite{kraus2012}, and each gap carries an integer label $q$ that is
the Chern number of that parent. The primes possess none of this structure:
their diffraction is closer to limit-periodic than to quasiperiodic, with no
phason, no synthetic dimension, and no gap label. A prime-indexed modulation
would open a dense family of weak gaps carrying no invariant. We note this
explicitly because the arithmetic appeal of such a construction is
considerable and its topological content is nil.

If the gaps of a Fibonacci isotope ribbon are genuinely labelled, their
positions are fixed by the sequence rather than by any realisation of
disorder, and edge branches must traverse gap $q$ exactly $q$ times per
phason cycle~\cite{kraus2012,apigo2019}. Section~\ref{sec:topo} tests both
statements. Section~\ref{sec:survival} first establishes whether such a
ribbon can exist at all.

\section{The survival window}
\label{sec:survival}

\subsection{Transfer matrix and localisation length}

For a chain of masses $M_n$ with nearest-neighbour force constant $K$ and
spacing $a$,
\begin{equation}
\begin{pmatrix} u_{n+1} \\ u_n \end{pmatrix} = T_n
\begin{pmatrix} u_n \\ u_{n-1} \end{pmatrix}, \quad
T_n = \begin{pmatrix} 2 - M_n\omega^2/K & -1 \\ 1 & 0 \end{pmatrix},
\end{equation}
the inverse localisation length is the Lyapunov exponent
$\gamma(\omega) = \lim_{N\to\infty} N^{-1}\ln\|\prod_n T_n\|$, with
$\xi = a/\gamma$ and in-gap transmission
$\mathcal{T}(L) \sim e^{-2L/\xi}$~\cite{anderson1958}. Force constants are
isotope-independent within the Born--Oppenheimer approximation, so a single
primitive-cell calculation serves every arrangement; only the diagonal mass
matrix is rebuilt.

\subsection{Gap strength from the mass contrast}

The $^{10}$BN and $^{11}$BN formula units (24.01 vs.\ 25.01~amu) give a
density contrast $\Delta\rho/\rho = 4.1\%$, an impedance contrast
$\Delta Z/Z = 2.04\%$, and single-interface reflection $|r| = 0.0102$
\cite{vuong2018}. Weak-reflection theory predicts
\begin{equation}
\xi_{\rm gap} = \frac{d}{2|r|}, \qquad
\frac{\Delta\omega}{\omega} = \frac{2|r|}{\pi} = 0.67\%,
\label{eq:gapstrength}
\end{equation}
with $d$ the modulation period; direct Lyapunov calculation on
$2\times10^5$-site chains reproduces both to $1$--$6\%$.

\subsection{The cost of topology}

Quasiperiodicity distributes reflection strength across a Cantor set of gaps
instead of one Bragg resonance~\cite{kohmoto1983,merlin1985}. A
same-frequency comparison settles its price: a periodic superlattice with its
Bragg gap at $0.25$~THz ($d_{\rm per} = 40.1$~nm) gives
$\xi_{\rm per} = 1.96\,\mu$m, matching the analytic $d/2|r|$ to three digits,
while the Fibonacci $(0,1)$ gap placed at the same frequency gives
$\xi_{\rm gap} = 4.31\,\mu$m. Topology costs
\begin{equation}
\xi_{\rm fib}/\xi_{\rm per} = 2.20
\end{equation}
at matched frequency and mass contrast.

\subsection{Two perturbations, opposite symmetries}

The isotope modulation spans the full ribbon cross-section and is strictly
diagonal in the transverse basis,
$V^{\rm iso}_{mm'} = -\omega^2\epsilon(z)\delta_{mm'}$: it opens a gap in
every channel independently, at a channel-dependent frequency --- the
smearing mechanism restated. Edge roughness breaks transverse translation
and is off-diagonal; its consequence depends qualitatively on the mode
count. In multi-mode operation it scatters between subbands --- genuine
decoherence. In single-mode operation there is nowhere to scatter but
backwards, and the mechanism reduces to one-dimensional Anderson
localisation, which \emph{adds} to $\gamma$ rather than destroying the gap.
A single-mode ribbon converts the dominant loss into one of a different
kind; establishing strict single-mode operation is accordingly the central
technical requirement.

\subsection{The tension and the working point}
\label{sec:tension}

Single-mode operation with free (Neumann) edges requires $f < v/2W$.
Coherence requires $\xi_{\rm dis} > 3\,\xi_{\rm gap}$. With
$\xi_{\rm gap} = C_\xi\lambda$ and the calibrated law of
Sec.~\ref{sec:calib}, the two combine into
\begin{equation}
\lambda^3 > 24\pi^2 C_\xi\, S_w(2q) \;\Rightarrow\;
\lambda \gtrsim 54\,[\Delta_r/\mathrm{nm}]^{2/3}~\mathrm{nm},
\end{equation}
a far softer wall than the linear bound the Ziman form would imply; the
change of exponent from $\Delta_r$ to $\Delta_r^{2/3}$ follows from the
$W^2$ law below.

The gaps of a Fibonacci chain sit on the Fourier module,
\begin{equation}
f_{(p,q)} = \frac{v}{2d}\,\bigl|p + q\,\tau^{-1}\bigr|,
\label{eq:module}
\end{equation}
$\tau = (1+\sqrt5)/2$. Placing the $q = 1$ gap at the target frequency fixes
$d = \tau^{-1} v/2f = 24.8$~nm ($99$ lattice sites per letter); a
module-targeted scan confirms all four principal gaps within $0.3\%$ of
Eq.~\eqref{eq:module}, with labels verified through the integrated density of
states (IDOS). Two findings deserve record. The $(1,0)$ gap --- the only one
periodic intuition would seek, at $v/2d$ --- is thirty times weaker than
$(0,1)$ at this decoration and carries $q = 0$: not topological, and
fortunately nearly absent. The $(1,-1)$ gap at $0.154$~THz is as strong as
$(0,1)$ ($\xi = 4.18\,\mu$m) with the opposite label; both lie below the
single-mode ceiling of a $30$~nm ribbon, so one device hosts two topological
gaps of opposite Chern number --- a built-in control for
Sec.~\ref{sec:expt}.

\begin{table}[b]
\caption{\label{tab:wp}Working point. $\xi_{\rm dis}$ from the calibrated
law of Sec.~\ref{sec:calib}; margins are $\xi_{\rm dis}/\xi_{\rm gap}$.}
\begin{ruledtabular}
\begin{tabular}{lc}
frequency / operating gap & $0.250$~THz, $(0,1)$, $q = 1$\\
ribbon width $W$ / period $d$ & $30$~nm / $24.8$~nm\\
$\xi_{\rm gap}$ & $4.31\,\mu$m\\
$\xi_{\rm dis}$ at $\Delta_r = 0.6$ ($1.0$)~nm & $76$ ($30$)\,$\mu$m\\
coherence margin & $18\times$ ($7.0\times$)\\
\end{tabular}
\end{ruledtabular}
\end{table}

\subsection{Calibration: breakdown of the ray formula}
\label{sec:calib}

The solver is validated three ways before any scan: it reproduces the
clean-ribbon limit exactly ($\mathcal{T} = N_{\rm prop}$ to five digits at
three grid spacings, every transmission eigenvalue equal to unity); it passes
a realization-matched grid-convergence gate ($1.3\%$ between $h = 1$ and
$0.5$~nm); and it reproduces the analytic single-bump Born reflection
$|r|^2 = (q^2/W^2)|\tilde a(2q)|^2$, converging with the grid to a uniform
$+3.7\%$ excess --- a genuine second-order correction whose role below is
decisive.

The definitive scan --- fifteen single-mode-locked parameter points, $24$
disorder seeds and six lengths each, every length window spanning
$(0.4\text{--}3)\,\xi$ of the law itself so that the decay is sampled rather
than extrapolated --- gives
\begin{align}
\frac{{\rm d}\ln\xi}{{\rm d}\ln\Delta_r} &= -1.91 \pm 0.06, &
\frac{{\rm d}\ln\xi}{{\rm d}\ln W} &= +1.87 \pm 0.08,\nonumber\\
\frac{{\rm d}\ln\xi}{{\rm d}\ln f} &= -1.384 \pm 0.095,
\end{align}
against the law's $(-2,\,+2,\,-2+4q^2L_c^2 = -1.38)$: all three within
tolerance, the frequency exponent to one part in a thousand. The absolute
magnitude matches with median $\xi_{\rm meas}/\xi_{\rm law} = 0.965$, while
the single-bump test independently predicts $1/1.037 = 0.964$ for the same
second-order excess: two measurements of one correction, one deterministic
and one ensemble-averaged, agreeing to $10^{-3}$. Across the scan the Ziman
formula underestimates $\xi$ by a median factor of $10.8$: in the single-mode
wave limit $\lambda \gg W$ a rough edge is Born-transparent, and the
ray-limit picture of diffuse wall collisions does not apply. The scaling
structure --- the rate $q^2 S_w(2q)/W^2$ --- agrees with the perturbative
treatment of roughness scattering in few-mode elastic waveguides given by
Santamore and Cross~\cite{santamore2001} in the context of the
thermal-conductance quantum, there for clamped elastic boundaries; the
closed zero-parameter form of Eq.~\eqref{eq:law}, its absolute
normalization, and its verification deep into the localized regime are the
additions made here. The residuals
retain a mild trend in $qL_c$ (from $0.78$ at $qL_c = 0.16$ to $1.12$ at
$0.42$), the expected fingerprint of second-order structure beyond the
leading law.

Methodological record --- three failure modes, each caught by a different
gate and none by inspection. A staircase-masked solver produced
grid-divergent spurious \emph{excess} scattering ($\xi = 120/42/16$~nm at
$h = 2/1/0.5$~nm; no continuum limit), caught by grid-convergence testing. A
naive finite-difference discretization of the transformed operator was
non-Hermitian and current-violating (transmission eigenvalues up to $1.26$)
while remaining grid-\emph{converged} and Born-consistent in total
transmission: invisible to convergence testing, exposed only by a unitarity
gate on the transmission-eigenvalue spectrum, and cured by assembling the
discrete operator from the quadratic form of the problem, which makes
Hermiticity ($R_{m+1} = P_m^{\mathsf T}$ by construction) and the exact
tilted-wall Neumann condition properties of the construction rather than of
a hand-derived stencil. Length windows tied to the Ziman estimate --- eleven
times too short --- sample only shallow decay and biased the $W$ exponent by
$+0.3$ in a preliminary pass: windows must track the law being tested. And a
scan whose ladders crossed a transverse-mode threshold returned exponents of
opposite sign on one axis ($+5.45$ and $-2.91$), the diagnostic signature of
an unenforced single-mode condition. 

The consequence is the central practical result of this section: coherence no
longer demands $\Delta_r \lesssim 0.13$~nm --- a single atomic row --- but
survives $\Delta_r = 0.6$~nm with an eighteen-fold margin and clears the
$3\times$ threshold at $1.0$~nm (Table~\ref{tab:wp},
Fig.~\ref{fig:survival}): within routine electron-beam lithography.

A scope statement, owed to the reader before the topology of
Sec.~\ref{sec:topo} is built on these numbers: everything above treats a
single scalar polarization with linear dispersion $v = 20$~nm/ps. A real
suspended h-BN ribbon carries four low-frequency branches --- longitudinal,
in-plane transverse, torsional, and a flexural branch whose dispersion is
quadratic and which dominates the low-frequency density of states. The two
mechanisms this paper rests on are polarization-generic: gap labelling
requires only one-dimensional quasiperiodicity of the mass sequence, and
Born transparency only $q$-selective backscattering from a weak boundary
perturbation, so the structure of Eq.~\eqref{eq:law} and of the module
\eqref{eq:module} carries over branch by branch. The numbers do not: mode
mixing among branches, the quadratic flexural dispersion, and
branch-dependent edge coupling will shift working-point frequencies,
$\xi$ values and margins, and the scalar model is best read as a
quantitative treatment of the longitudinal channel and a structural
template for the rest. A full elastic calculation is the natural next step
and lies outside this paper.

\section{Topological verification}
\label{sec:topo}

\subsection{Gap labelling, falsifiably}

Within each gap the IDOS is constant and restricted by the gap-labelling
theorem to
$\mathcal{N} = (p + q\,\tau^{-1}) \bmod 1$~\cite{bellissard1992}. We compute
$\mathcal{N}$ by Sturm node counting on the transfer-matrix recursion
($\mathcal{O}(N)$, no diagonalisation); Table~\ref{tab:gaps} lists the
twelve strongest gaps of a $6\times10^4$-site chain
(Fig.~\ref{fig:gamma}). The two strongest take the textbook values
$\tau^{-1}$ and $1 - \tau^{-1}$.

\begin{figure}[t]
\includegraphics[width=\columnwidth]{fig_gamma_idos.pdf}
\caption{\label{fig:gamma}Lyapunov spectrum $\gamma(\omega)$ (top) and IDOS
staircase (bottom) of the Fibonacci chain, with the twelve labelled gaps of
Table~\ref{tab:gaps}. Inset: gap-labelling residuals, measured gaps against
$3\times10^3$ random controls (log axis) --- the falsifiability argument of
Sec.~\ref{sec:topo} in one picture.}
\end{figure}


Taken at face value the test is vacuous: $\mathbb{Z} + \tau^{-1}\mathbb{Z}$
is dense in $[0,1]$, so some pair fits any number. The content is
quantitative. Against $5\times10^3$ uniform controls under the identical
fitting procedure ($|q| \le 30$): measured gaps show median residual
$2.95\times10^{-6}$ (maximum $5.95\times10^{-6}$); random controls,
$4.08\times10^{-3}$ --- a separation of $1385$; the probability that a
random value fits as well as the \emph{worst} measured gap is
$p = 1.4\times10^{-3}$. We recommend this control whenever gap labelling is
claimed numerically.

\begin{table}[t]
\caption{\label{tab:gaps}The twelve principal gaps ($N = 6\times10^4$
sites). $\xi$ at gap centre; residual is
$|\mathcal{N} - (p + q\tau^{-1})|$ for the best pair with $|q| \le 30$.}
\begin{ruledtabular}
\begin{tabular}{cccccc}
$f$ (THz) & $\gamma$ (site$^{-1}$) & $\xi$ (nm) & IDOS & $(p,q)$ & resid.\\
\hline
14.386 & 0.00839 & 29.9 & 0.38197 & $(1,-1)$ & $7{\times}10^{-7}$\\
17.213 & 0.00290 & 86.5 & 0.47213 & $(-2,4)$ & $3{\times}10^{-6}$\\
18.792 & 0.00356 & 70.3 & 0.52787 & $(3,-4)$ & $3{\times}10^{-6}$\\
19.693 & 0.00170 & 147.4 & 0.56230 & $(-5,9)$ & $6{\times}10^{-6}$\\
21.040 & 0.01744 & 14.4 & 0.61803 & $(0,1)$ & $7{\times}10^{-7}$\\
22.222 & 0.00284 & 88.1 & 0.67377 & $(5,-7)$ & $5{\times}10^{-6}$\\
22.859 & 0.00345 & 72.6 & 0.70820 & $(-3,6)$ & $4{\times}10^{-6}$\\
23.760 & 0.01155 & 21.7 & 0.76393 & $(2,-2)$ & $1{\times}10^{-6}$\\
24.222 & 0.00185 & 135.1 & 0.79837 & $(-6,11)$ & $7{\times}10^{-6}$\\
24.479 & 0.00239 & 104.7 & 0.81967 & $(7,-10)$ & $7{\times}10^{-6}$\\
24.826 & 0.00812 & 30.8 & 0.85410 & $(-1,3)$ & $2{\times}10^{-6}$\\
25.239 & 0.00516 & 48.5 & 0.90983 & $(4,-5)$ & $3{\times}10^{-6}$\\
\end{tabular}
\end{ruledtabular}
\end{table}

\subsection{Phason pumping: the full counter-propagating structure}
\label{sec:pump}

The phason angle $\phi$ of the cut-and-project word
$s_n(\phi) = \lfloor (n{+}2)\tau^{-1}{+}\phi \rfloor -
\lfloor (n{+}1)\tau^{-1}{+}\phi \rfloor$ is the synthetic dimension; a full
cycle traverses the hull, and edge branches must cross gap $q$ exactly $q$
times~\cite{kraus2012}. Tracking such branches requires care: sorting by
frequency exchanges identities at band crossings, so we continue states by
eigenvector overlap with Hungarian assignment, terminating any branch whose
best overlap falls below $0.15$, and record left- and right-end weights
separately (Fig.~\ref{fig:phason}).

\begin{figure}[t]
\includegraphics[width=\columnwidth]{fig_phason.pdf}
\caption{\label{fig:phason}Phason spectrum of the $(0,1)$ gap: frequency
against $\phi$, bulk states as points, edge-bound branches coloured by the
\emph{instantaneous} dominant end (blue left, orange right), line width
$\propto$ edge weight. The colour change at the in-gap degeneracy is the
identity exchange of Sec.~\ref{sec:topo}C.}
\end{figure}


A first pass on the $(0,1)$ gap ($610$-site free chain, $400$ steps) tracked
$27$ branches and reported a single edge-bound branch --- right end, weight
$0.721$, sweeping $86\%$ of the gap with one kink --- stable across chain
lengths $610$--$1597$ and steps $400$--$800$. Taken alone this suggested a
one-sided pump. Two basis-independent measurements correct that reading.
\emph{Spectral flow}: counting eigenvalues below each gap centre at $600$
phason steps, every one of the twelve gaps shows exactly $|q|$ downward and
$|q|$ upward crossings per cycle --- $1{+}1$ for $(0,1)$, $3{+}3$ for
$(-1,3)$, up to $11{+}11$ for $(-6,11)$ --- without exception.
\emph{Crossing-moment classification}: diagonalizing with eigenvectors at
each detected crossing, the descending crossings of every $q > 0$ gap are
right-bound and the ascending left-bound, exactly mirrored for $q < 0$; for
$(-1,3)$ all six crossings alternate side-perfectly, and the $(0,1)$ pair
carries weights $0.72$ and $0.71$. The counter-propagating pair exists in
full; the per-edge spectral flow equals $\pm q$; the bulk--boundary
correspondence holds in its standard form for all twelve labels.

\subsection{A cautionary note on overlap tracking}

Why did the first pass see one branch? The two counter-propagating edge
modes hybridize where they are degenerate inside the gap; at that avoided
crossing, overlap continuation --- by construction --- follows the smoothly
deforming hybrids and exchanges identities, stitching the descending halves
of the two branches into a single apparent branch with one kink: precisely
the phenomenology reported. The spectral-flow classification never asks for
continuity between steps and is immune. Under fixed (Dirichlet) ends the
pair persists (crossings at $\phi = 0.833, 0.932$ against $0.618, 0.147$
free), so termination relocates the crossings but removes nothing. We record
the artefact because Hungarian-overlap continuation is standard practice for
pump spectra, and it fails silently in exactly the situation ---
counter-propagating branches meeting inside a gap --- that topological pumps
guarantee.

\section{Experimental design}
\label{sec:expt}

\subsection{Frequency band; exclusion of transient gratings}

The working gap sits at $0.250$~THz, $\lambda = 80$~nm. The natural
wavevector-selective probe, the transient thermal grating, is excluded by
optics: two beams of wavelength $\lambda_{\rm opt}$ cannot write a grating
finer than $\lambda_{\rm opt}/2$, so a $400$~nm pump bottoms out near
$200$~nm --- $2.5\times$ too coarse --- and $266$~nm reaches only $133$~nm.
No geometry closes the factor. The measurement must be performed in the time
domain: asynchronous optical sampling (ASOPS) between two femtosecond
oscillators offset by a few kHz~\cite{bartels2007} provides
$\sim\!100$~fs effective resolution across the full $12.5$~ns repetition
window, sampling a $4$~ps acoustic period over $\sim\!3\times10^3$ cycles
--- spectral resolution $\sim\!0.3$~GHz.

\subsection{Cavity-mode spectroscopy of the gap}

Aluminium pads ($\sim$20~nm) on both ends serve as broadband phonon
sources~\cite{thomsen1986}, detection mirrors, and the reflectors of an
acoustic Fabry--P\'erot cavity; the ribbon interior is interrogated through
its mode spectrum alone. A cavity of $L = 20\,\mu$m has free spectral range
$v/2L = 0.5$~GHz; the $(0,1)$ gap ($1.7$~GHz) expels three to four
consecutive modes, and with $\xi_{\rm gap} = 4.31\,\mu$m the in-gap
transmission $e^{-2L/\xi} \approx 10^{-4}$: modes are removed, not shifted.
The observable is differential, as in Ref.~\cite{paper1}: the Fibonacci
ribbon is paired on-chip with a composition-matched random permutation of
the same isotope content; boundary, transducer and anharmonic contributions
cancel to first order (Fig.~\ref{fig:expt}). The resolution budget is
stated plainly: five elements across the $(0,1)$ gap, three across
$(1,-1)$ --- sufficient to establish absence of modes, insufficient for
lineshape work.

\begin{figure}[t]
\includegraphics[width=\columnwidth]{fig_experiment.pdf}
\caption{\label{fig:expt}Proposed measurement (schematic). (a)~Ribbon with
Fibonacci isotope letters and Al transducers. (b)~Illustrative cavity
spectra: composition-matched random control (grey, full mode comb) against
Fibonacci (blue, gap emptied), with the $\phi$-dependent edge mode inside
the gap.}
\end{figure}


\subsection{Edge-state pumping by a phason array}

The phason angle is fixed lithographically, one value per ribbon; an array
of $\sim$20 ribbons at $\phi = 0, 0.05, \ldots, 0.95$ converts the pump
cycle into chip position. Section~\ref{sec:pump} predicts an in-gap
resonance whose frequency advances monotonically across the array, present
for most of the cycle. Two structural controls come free. The $(1,-1)$ gap
in the same ribbons carries the opposite label, so its in-gap resonance must
wind oppositely across the same array; any trivial systematic --- strain,
thickness drift, composition gradient --- moves both gaps identically and
is excluded. And the counter-propagating structure is directly testable:
the descending resonance ($\phi \approx 0.62$ for $(0,1)$) lives at one end
and the ascending one ($\phi \approx 0.15$) at the other, so exchanging
pump and probe transducers must swap which portion of the array shows the
stronger in-gap response --- a signature no single-sided mechanism
reproduces.

\subsection{Sample specification}

$W = 30$~nm (single-mode to $0.267$~THz); $L = 20\,\mu$m
($4.6\,\xi_{\rm gap}$); $d = 24.8$~nm ($99$ sites per letter);
$\sim$800 letters per ribbon ($F_{16} = 987$ truncated);
$\Delta_r \le 0.6$~nm rms ($18\times$ margin); $T \le 4$~K. The isotope
pattern is one-dimensional and in-plane, within electron-beam masking
resolution, with sequential two-precursor regrowth as the candidate route.
The contrast with Ref.~\cite{paper1} bears repeating: that design required
$2.4\times10^5$ independently specified units over $300\,\mu$m; this one,
$\sim$800 letters over $20\,\mu$m at mesoscopic pitch, with edge tolerances
relaxed an order of magnitude by the physics of Sec.~\ref{sec:calib}. The
tension that forced the working point down in frequency also carried it into
fabricability.

\begin{acknowledgments}
Numerical calculations were performed with open-source software. All code,
calibration data, and figure-generation scripts sufficient to reproduce every
number in this manuscript are available at
\url{https://github.com/Ruqing1963/hbn-fibonacci-topological-phonons}.
\end{acknowledgments}

\bibliography{refs}

\end{document}
